\scanpage{27}
Le fait qu'il soit gradué est immédiat. Il contient le groupe des cycles lin. équiv. à $0$, comme il résulte aussitôt des définitions, et du fait que sur $T=\mathbb P^1$, tout $0$-cycle $\nu$ de degré $0$ satisfait $S_T(\nu)=0$. \deleted{Ainsi $N(X)$ est bien un \unsure{k-groupe|s-groupe} de $X$, contenant}

Pour prouver (iii), notons $z$ un cycle de la classe envisagée dans $A^*(Y)$, et considérons un $Z$ sur $Y\times T$ et un $\nu$ comme dans 1.1., donc $z=Z(\nu)$. Si nous changeons $Z,\nu$ dans leur classe d'équivalence linéaire \inserted{en $Z',\nu'$}, $Z'(\nu')$ sera linéairement équivalent à $z$, et le Pb est le $=$ pr $z$ et $z'$. Il suffit maintenant de prouver \deleted{\unsure{et}} qu'il existe $Z',\nu'$ tels que $(f\times\operatorname{id}_T)^*(Z')=Z'_1$ et $Z'(\nu')$, \deleted{$Z'_1(\nu')$} soient définis; en effet, en vertu de la théorie des intersections, on aura alors
\[
f^*(Z'(\nu'))=\deleted{\unsure{f^*}}Z_1(\nu'),
\]
or le premier \unsure{mb.} \deleted{\unsure{représente est dans le}} représente $f^*(z)$ dans $A^*(X)$, \deleted{le deuxième est}

\margin{\diagram{\begin{tikzcd}[ampersand replacement=\&]
X\arrow[r,dash,"{\unsure{-|\to}}" description]\&Y
\end{tikzcd}}}
\margin{\diagram{\begin{tikzcd}[ampersand replacement=\&]
X\times T\arrow[rr,Leftrightarrow,"{\unsure{\Leftrightarrow|\Rightarrow}}" description]\arrow[dr]\& \&Y\times T\arrow[dl]\\
\&\deleted{\unsure{T}}T\&
\end{tikzcd}}}

\mspage{26}
